% !TEX root = ../main.tex
% =====================================================================
% CAPITOLO 3 — LA SUITE DI AGENTI VERTICALI
% =====================================================================

\chapter{La suite di agenti verticali}
\label{chap:agenti}

Questo capitolo descrive la prima fase del progetto: la progettazione e la
realizzazione della suite di agenti verticali. Il percorso segue l'ordine
delle decisioni effettivamente prese: per primi i principi guida che hanno
condizionato ogni scelta successiva e l'architettura d'insieme che ne
deriva; quindi le fondazioni trasversali comuni a tutti i servizi (stack
tecnologico, organizzazione del codice, infrastruttura come codice, identità
e sicurezza) e infine il dettaglio di ciascuno dei quattro agenti. 
Chiudono il capitolo il
pannello di amministrazione e le alternative progettuali considerate e
scartate.

% =====================================================================
\section{Principi guida}
\label{sec:agenti-principi}
% =====================================================================
\label{subsec:agenti-principi-guida}

La progettazione è stata guidata da una serie di fondamenti trasversali,
stabiliti prima di avviare lo sviluppo e mantenuti durante le fasi successive.

\begin{description}

  \item[Rispetto dei principi etici e cooperativi]
    Gli agenti LLM sono intrinsecamente portati a esprimere bias (Sez.~\ref{sec:bg-llm}).
    Sviluppando agenti interni, il  "comportamento" (definito da guardrail 
    e prompt di sistema) si può definite collegialmente fra gli uffici interni (Etica, Sicurezza IT),
    integrando direttive esplicite che allineano ogni agente all'identità e ai valori
    della cooperativa e in conformità con i requisiti di trasparenza e \emph{human oversight}
    dell'European AI Act~\cite{eu2024aiact}, con cicli di aggiornamento rapidi e controllati.
    Inoltre, offrire un pacchetto di assistenti interni mitiga questo previene il
    ricorso incontrollato a strumenti AI esterni non autorizzati (\emph{Shadow AI}), riducendo i rischi per
    la sicurezza del dato.

  \item[Governance facilitata e efficace]
    Le configurazioni del prompt di sistema (la ``personalità'' dell'agente) e i controlli su
    utilizzi impropri o sensibili devono essere configurabili da una dashboard apposita, 
    usabile anche dagli uffici preposti privi di competenza tecnica, che consenta inoltre di monitorare
    gli eventi segnalati come ``pericolosi''. Tutto il codice sorgente risiede nel repository 
    Git della Cooperativa ed è di sua proprietà.

  \item[Implementazione agnostica e sostenibile]
    Le soluzioni AI ``out of the box'' dei fornitori commerciali comportano spesso \emph{vendor lock-in}
    (difficile portabilità verso altre infrastrutture) e costi fissi di sottoscrizione anche con
    utilizzo scarso. La soluzione realizzata vuole invece essere portabile verso qualunque
    architettura o fornitore di LLM e consentire un controllo dinamico dei costi: un budget di utilizzo (``token'')
    per utente, definito dagli amministratori e modificabile alla bisogna, e ovunque possibile un
    \emph{routing} che indirizzi le richieste a bassa complessità verso modelli più economici.

  \item[Sovranità del dato]
    Nessun dato deve lasciare il perimetro cloud europeo, in conformità al Regolamento Generale sulla Protezione dei Dati (GDPR).

  \item[Modularità verticale]
    Ogni dominio informativo o funzionale deve essere affidato a un agente autonomo e indipendente,
    con il proprio meccanismo di grounding, il proprio ciclo di vita applicativo e la propria infrastruttura di deployment:
    in modo che sia possibile effettuare \emph{prompt engineering} e \emph{few-shot learning}~\cite{brown2020gpt3}
    ad hoc per il singolo caso d'uso.
    I sistemi devono essere il più possibile isolati per consentire di aggiornare, monitorare e
    scalare ogni servizio in modo indipendente.

  \item[Automazione dei rilasci]
    Le evolutive o correzioni di bug devono potere essere rilasciate in maniera trasparente verso l'utente 
    senza causare disservizio,
    in maniera automatizzata per ridurre al minimo la possibilità di errore umano,
    e con uno step approvativo per il passaggio in produzione.

  \item[Accesso privato e con controllo centralizzato]
    L'intero ecosistema deve potere essere accessibile solo da rete interna o VPN di Coop Alleanza:
    l'infrastruttura utilizzata deve quindi evitare l'esposizione su rete pubblica.
    Il controllo dell'abilitazione degli utenti all'uso degli assistenti deve integrarsi con il sistema di Identity Management
    già esistente nell'organizzazione, basato sui protocolli OAuth~2.0~\cite{rfc6749oauth2} e OpenID Connect~\cite{oidc1_0core}
    su Microsoft Entra ID~\cite{msft_entra_id}.

\end{description}

% =====================================================================
\section{Architettura d'insieme}
\label{subsec:agenti-architettura-insieme}
% =====================================================================

L'ecosistema è accessibile da un unico ``portale'', per consentire modularità: 
eventuali nuovi agenti si innestano nel \emph{framework} esistente, così
che l'utente finale possa accedervi in maniera pratica e intuitiva.

Pur restando centrale il principio di agnosticità e portabilità, tutte le implementazioni usano
la piattaforma Cloud di Google~\cite{gcp_vertexai,gcp_cloudrun}, al momento l'unico fornitore
attivo in Coop Alleanza e adeguato alle necessità.

\begin{figure}[htbp]
  \centering
  \includegraphics[width=0.9\textwidth]{fig1}
  \caption[Architettura d'insieme della suite di agenti]{Architettura
    della suite di agenti verticali: portal con autenticazione federata,
    istanze serverless indipendenti e infrastruttura
    condivisa per database, data warehouse e servizi AI.}
  \label{fig:architettura_suite}
\end{figure}

% =====================================================================
\section{Fondazioni trasversali: stack, infrastruttura, sicurezza}
\label{sec:agenti-stack}
% =====================================================================

Prima di esaminare i singoli agenti, questa sezione descrive gli elementi
trasversali del progetto,presentati qui una sola volta: l
o stack tecnologico, l'organizzazione del
codice, l'infrastruttura come codice con la pipeline di rilascio, e il
modello di identità e autorizzazione.

\subsection{Stack tecnologico}
\label{subsec:agenti-stack-tech}

L'implementazione dell'ecosistema di agenti è basata su uno stack cloud-native
 omogeneo per minimizzare la diversità fra servizi (riducendo la complessità
per quanto possibile) e per delegare ai \emph{managed services} di Google Cloud Platform tutti i compiti
infrastrutturali.

\begin{itemize}
  \item \textbf{Backend.} Tutti gli agenti sono scritti in Python~3.10+ ed
        esposti tramite il framework asincrono FastAPI~\cite{fastapi_docs}. 
        La scelta è motivata
        dalla natura prevalentemente \emph{I/O bound} delle interazioni con i
        Large Language Model: la generazione di una risposta richiede da
        centinaia di millisecondi a decine di secondi, quasi interamente spesi
        in attesa del modello, dei motori di retrieval o di BigQuery. Il modello
        di concorrenza \texttt{async/await} consente di sostenere centinaia di
        richieste concorrenti per istanza Cloud Run senza ricorrere a
        thread o processi multipli. L'accesso al database relazionale avviene
        in modo asincrono tramite SQLAlchemy~2 e il driver
        \texttt{asyncpg}, con \emph{pool} pre-pingati e ricicli automatici.
  \item \textbf{Frontend.} Le interfacce sono React~19 con TypeScript,
        impacchettate con Vite e stilizzate con Tailwind~CSS. La scelta
        di un \emph{bundler} ESM moderno è funzionale alla pipeline CI/CD:
        il build è di pochi secondi, riducendo i tempi di rilascio.
  \item \textbf{Database e Data Warehouse.} Cloud SQL per
        PostgreSQL~15~\cite{gcp_cloudsql} è il \emph{system of record}
        applicativo: vi risiedono utenti, conversazioni, messaggi, quote,
        ACL, audit log e parametri di configurazione. Google
        BigQuery~\cite{gcp_bigquery} è invece il database analitico
        utilizzato dagli agenti Text-to-SQL (AllyCare e SysAid), con
        dataset alimentati in modo continuo da flussi appositi.
  \item \textbf{Large Language Models.} L'implementazione poggia su due
        famiglie di modelli erogati come \emph{managed APIs}: Gemini~2.5
        Flash e Pro~\cite{team2024gemini} su Vertex AI~\cite{gcp_vertexai},
        e Claude Sonnet~4.6~\cite{anthropic2024claude}
        attraverso l'SDK Anthropic veicolato dal canale
        \texttt{AnthropicVertex}~\cite{anthropic_docs}. La coesistenza
        dei due fornitori in SysAid (Sez.~\ref{sec:agenti-sysaid-arch})
        concretizza un caso di agnosticità come richiesto dai principi
        guida (Sez.~\ref{sec:agenti-principi}).
  \item \textbf{Sviluppo di agenti.} Viene utilizzato il \emph{Google Agent
        Development Kit}~(ADK)~\cite{google_adk_docs}, una libreria aperta,
        ufficiale e portabile che astrae la gestione della sessione, l'invocazione
        ciclica del modello e l'integrazione di \emph{tool} basati su
        \emph{function calling}.
 \item \textbf{Infrastruttura.} Ogni componente applicativo è un servizio
        serverless eseguito su Google Cloud Run~\cite{gcp_cloudrun},
        descritto interamente da codice Terraform~\cite{terraform_docs}.
\end{itemize}

\subsection{Organizzazione del codice}
\label{subsec:agenti-multirepo}

Il codice sorgente dell'intera suite è organizzato in repository Git distinti
(Fig.~\ref{fig:multirepo_tree}), ciascuno autonomo e con il proprio ciclo di
vita, e raccolti in un unico progetto Bitbucket (Fig.~\ref{fig:progetto_bitbucket}). Ogni repository
contiene il proprio \texttt{Dockerfile}, le proprie dipendenze e il proprio
file di descrizione della pipeline (\texttt{cloudbuild.yaml}),
ed è collegato a un trigger Cloud Build dedicato. Questa separazione consente
rilasci indipendenti: un aggiornamento a CooPolicy non richiede il rebuild
di Vera, e un hotfix su SysAid può essere rilasciato senza coordinamento
con il frontend del portale. I branch \texttt{test} e \texttt{main} di ogni
repository corrispondono rispettivamente all'ambiente di collaudo e a quello
di produzione; il branch develop può invece essere sfruttato per i test in locale.

\begin{figure}[h]
  \centering
  \begin{lstlisting}[language=bash]
vera-backend/              # Vera Backend (FastAPI + ADK + Gemini)
vera-frontend/             # Vera Frontend (React)
vera-admin-backend/        # Pannello amministrativo di Vera Backend
vera-admin-frontend/       # Pannello amministrativo di Vera Frontend
sysaid-backend/            # SysAid Backend unificato (Gemini + Sonnet)
sysaid-frontend/           # SysAid Frontend
coopolicy-backend/         # CooPolicy Backend (FastAPI + RAG)
coopolicy-frontend/        # CooPolicy Frontend
allycare-backend/          # AllyCare Backend (FastAPI + ADK + BigQuery)
allycare-frontend/         # AllyCare Frontend
portal-frontend/           # Portale utenti (MSAL)
admin-portal-frontend/     # Portale amministrativo
infrastructure/            # Terraform per infrastruttura
  \end{lstlisting}
  \caption[Elenco dei repository della suite (agenti)]{Elenco dei repository della suite nella ``fase'' degli agenti verticali. Ogni servizio è un progetto Git autonomo con proprio ciclo di CI/CD e branch per ambiente (\texttt{test}, \texttt{main}). I repository dell'orchestratore sono inseriti dal Capitolo~\ref{chap:orion}.}
  \label{fig:multirepo_tree}
\end{figure}

\begin{figure}[htbp]
  \centering
  \includegraphics[width=0.9\textwidth]{progetto-bitbucket}
  \caption[Progetto Bitbucket della suite]{Il progetto Bitbucket che raccoglie i repository della suite: ogni servizio è un repository autonomo con il proprio ciclo di CI/CD.}
  \label{fig:progetto_bitbucket}
\end{figure}

I backend più complessi (Vera e Vera-Admin) adottano una separazione interna del codice
più strutturata: la directory \texttt{api/} contiene esclusivamente i router
FastAPI; \texttt{services/} le funzioni di dominio prive di stato HTTP;
\texttt{models/} le entit\`a SQLAlchemy; \texttt{schemas/} i modelli Pydantic
per validazione di input e output; \texttt{core/} le utility trasversali
(database, dipendenze, sicurezza). I backend più piccoli (CooPolicy, AllyCare e
SysAid) condividono la stessa filosofia in forma compressa, con
un singolo \texttt{main.py} che monta i \emph{router} e alcuni moduli
ausiliari (\texttt{auth.py}, \texttt{tools.py}, \texttt{agent.py}, \texttt{quote.py}).

\subsection{Infrastruttura come codice e pipeline CI/CD}
\label{subsec:agenti-iac}

L'intera infrastruttura GCP è descritta in Terraform~\cite{terraform_docs} (Fig.~\ref{fig:arch_terraform}) ed è
parametrizzata su due ambienti, \texttt{test} e \texttt{prod}, gestiti come
progetti indipendenti con \emph{state} remoto su un bucket Cloud
Storage. Le risorse provisionate comprendono: l'abilitazione delle API GCP
necessarie, le istanze Cloud Run con un service account condiviso, l'istanza
Cloud SQL con il relativo \emph{Private Service Connect} consumer endpoint, le
politiche di Identity and Access Management (IAM), il bucket di Artifact
Registry per le immagini Docker, il template Model Armor, la rete (subnet,
indirizzi IP riservati, regole di forwarding) e il bilanciatore di carico HTTPS
interno con il certificato TLS recuperato da Secret Manager (Certificato fornito dall'IT di Coop Alleanza).

Le decisioni di sicurezza fondamentali sono codificate direttamente nel
modulo Terraform e dunque non sono modificabili senza una \emph{pull request}
revisionata:

\begin{enumerate}
  \item \textbf{Egress limitato a indirizzi privati.} Ogni servizio Cloud Run
        è collegato alla VPC interna mediante \emph{Direct VPC Egress} con
        \texttt{egress = PRIVATE\_RANGES\_ONLY}: il container può raggiungere
        gli endpoint privati GCP (Vertex AI, BigQuery, Cloud SQL, RAG Engine)
        ma non internet pubblico. Il vincolo è particolarmente importante per
        gli agenti Text-to-SQL e per Vera, dove un'eventuale estrazione di dati verso
        l'esterno via DNS o webhook è di fatto resa impossibile a livello di
        rete.
  \item \textbf{Database non esposto su IP pubblico.} L'istanza Cloud SQL ha
        \texttt{ipv4\_enabled = false}: la connessione passa solo attraverso
        \emph{Private Service Connect} (PSC), con un \emph{forwarding rule}
        regionale verso il \emph{service attachment} di Cloud SQL:
        è, in sostanza, anch'esso esposto in sola rete privata.
  \item \textbf{Bilanciatore di carico interno.} Il \emph{frontend} è
        accessibile solo via un \emph{Application Load Balancer regionale
        interno} (\texttt{INTERNAL\_MANAGED}). L'esposizione
        avviene perciò solo all'interno della rete aziendale o, per gli
        utenti remoti, dietro VPN, soddisfacendo il requisito di
        accesso privato (Sez.~\ref{sec:agenti-principi}).
  \item \textbf{Segregazione di credenziali.} Tutte le password generate
        (database, JWT) sono prodotte da \texttt{random\_password},
        depositate in Secret Manager e iniettate nei container via
        \texttt{value\_source.secret\_key\_ref}: nessun segreto
        attraversa mai il file system dello sviluppatore.
\end{enumerate}

\begin{figure}[htbp]
  \centering
  \includegraphics[width=0.9\textwidth]{arch-terraform}
  \caption[Architettura GCP descritta in Terraform]{Architettura completa dell'infrastruttura GCP della suite, così come descritta nei moduli Terraform: Cloud Run, Cloud SQL, Load Balancer interno, Secret Manager, Vertex AI.}
  \label{fig:arch_terraform}
\end{figure}

La pipeline CI/CD (Fig.~\ref{fig:cicd_pipeline}) si attiva ad ogni \emph{push}: in ambiente di test sul branch
\texttt{test}, in produzione sul branch \texttt{main}. Cloud Build, configurato
tramite un \emph{trigger} per ciascun servizio collegato al rispettivo
repository Bitbucket, esegue i passi descritti nel \texttt{cloudbuild.yaml}
locale: build dell'immagine Docker, push su Artifact Registry e aggiornamento
dell'immagine attiva del relativo servizio Cloud Run. Per l'ambiente di
produzione, il \emph{trigger} ha \texttt{approval\_required = true}, che
implementa lo step approvativo richiesto dal principio di
\emph{Automazione dei rilasci} (Sez.~\ref{sec:agenti-principi}).
Tale rilascio potrà essere poi ``sbloccato'' solo dagli amministratori di sistema della Cooperativa.

\begin{figure}[htbp]
  \centering
  \includegraphics[height=0.75\textheight,keepaspectratio]{gitflow}
  \caption[Branching strategy e pipeline CI/CD]{La branching strategy adottata: ogni push sul branch \texttt{test} avvia un deploy automatico sull'ambiente di collaudo; il passaggio in produzione avviene tramite merge su \texttt{main} e richiede l'approvazione esplicita di un revisore prima del deploy su Cloud Run.}
  \label{fig:cicd_pipeline}
\end{figure}

\subsection{Identità, autenticazione e autorizzazione}
\label{subsec:agenti-auth}

L'autenticazione è interamente delegata a Microsoft Entra ID, che è l'unico
\emph{tenant} di identità usato in Cooperativa.
La scelta del flusso e della libreria segue lo standard di mercato per le \emph{single-page
applications}: il browser ottiene un \emph{access token} via
OAuth~2.0~\cite{rfc6749oauth2} e OpenID Connect~\cite{oidc1_0core} con
\emph{Authorization Code Flow + PKCE}, gestito dalla
\emph{Microsoft Authentication Library} (MSAL)~\cite{msft_msal} caricata nel
portale e nel pannello di amministrazione. Il token è un JWT firmato in
RS256~\cite{rfc7519jwt} ed è il vettore di autorizzazione che attraversa la
suite.

Sul lato server, la verifica del token Entra ID è affidata alla funzione
\texttt{verify\_entra\_token} all'atto del login federato ed esegue le verifiche
minime ma non negoziabili: scarica e mette in cache
le chiavi pubbliche da \texttt{JWKS} esposto da Microsoft, identifica la
chiave corretta tramite il campo \texttt{kid} dell'header, decodifica il
token con la libreria \texttt{python-jose} validandone audience e
firma e, infine, valida manualmente l'\texttt{issuer} contro la coppia
\texttt{(sts.windows.net, login.microsoftonline.com)}: questa
ridondanza è necessaria perché Entra ID emette token con \texttt{iss}
v1 o v2 a seconda dell'\emph{endpoint} usato, e l'opzione \texttt{verify\_iss}
del decoder accetta un solo valore alla volta. Una rotazione di chiavi non
prevista forza un secondo \emph{fetch} del JWKS prima di restituire un
errore 401: la cache non è quindi un \emph{single point of failure}.
Superata la verifica, il backend emette un JWT di sessione applicativo, firmato
con un segreto interno, che le richieste successive presentano alla dipendenza
FastAPI montata su ogni endpoint protetto. Accanto al login federato, un meccanismo
di autenticazione locale con credenziali dedicate resta riservato all'utenza
amministrativa di servizio (inizializzazione e accesso di emergenza).

L'autorizzazione si basa su un modello semplificato: un unico gruppo
Active Directory (``AI Users'') controlla l'accesso all'intera piattaforma.
Il claim JWT \texttt{group AI Users} contiene il GUID del gruppo;
il portale verifica la presenza dell'utente nel gruppo e, in caso positivo,
rende visibili gli agenti. I backend replicano la stessa verifica
sui propri endpoint. I claim
\texttt{preferred\_username}, \texttt{Ruolo} e \texttt{Direzione},
propagati dai sistemi HR a Entra ID, sono utilizzati per personalizzare
le risposte degli agenti (es. tono, contesto organizzativo) ma non per
controllare l'accesso.

Le Fig.~\ref{fig:login1} e~\ref{fig:login2} mostrano rispettivamente il punto di accesso al portale, da cui parte il flusso OAuth~2.0 + PKCE, e la schermata di autenticazione Entra ID personalizzata con il brand aziendale.

\begin{figure}[htbp]
  \centering
  \includegraphics[width=0.75\textwidth,height=0.4\textheight,keepaspectratio]{login1}
  \caption[Portale agenti: pulsante di accesso]{Il portale agenti prima dell'autenticazione: il pulsante avvia il flusso OAuth~2.0 + PKCE delegato a Microsoft Entra ID tramite MSAL.}
  \label{fig:login1}
  \vspace{1em}
  
\includegraphics[width=0.75\textwidth,height=0.4\textheight,keepaspectratio]{login2}
  \caption[Schermata di autenticazione Microsoft personalizzata]{La schermata di autenticazione Microsoft Entra ID personalizzata con il brand dell'organizzazione:
  l'utente accede con le credenziali aziendali esistenti e controllo MFA, senza registrazione aggiuntiva.}
  \label{fig:login2}
\end{figure}

Il consumo di token è controllato da un servizio di quota su una \emph{tabella condivisa}
(\texttt{token\_budget}): ogni utente ha un budget mensile unico, indipendente dall'agente, su
fasce (\emph{Bronze}, \emph{Silver}, \emph{Gold}) modificabili dall'amministratore. Prima di ogni
chiamata il backend verifica il budget residuo e al termine contabilizza i token effettivi,
realizzando trasversalmente il principio di \emph{Implementazione agnostica e sostenibile}
(Sez.~\ref{sec:agenti-principi}).

Il prompt di sistema --- in ogni agente tranne CooPolicy, dove l'ottimizzazione avviene più a monte,
sul retrieval RAG --- impone una disciplina di output esplicita: ogni risposta apre con una sezione
\texttt{\#\#\# RAGIONAMENTO} (obiettivo, scelte, strategia), chiama i \emph{tool} appropriati e solo
dopo i risultati produce la sezione \texttt{\#\#\# RISPOSTA}. È una forma di Chain-of-Thought
esplicito~\cite{wei2022cot} con due tag che separano pianificazione e output; una funzione
\texttt{\_clean\_answer} collassa il blocco di ragionamento da quanto mostrato all'utente,
lasciandolo disponibile per l'auditing ma mantenendo pulita l'esperienza.

Le chiamate al modello sono racchiuse in un \texttt{try/except} che riconosce gli errori transitori
(\texttt{RESOURCE\_EXHAUSTED}, \texttt{429}, \texttt{503}, \emph{rate limit}) e li ripete con
\emph{backoff} lineare fino a tre tentativi, così che i picchi di traffico verso Vertex AI non
degradino il servizio percepito. La stessa strategia è usata da CooPolicy verso il RAG Engine e da
AllyCare e SysAid nelle query BigQuery.

% =====================================================================
\section{Il portale utente}
\label{sec:agenti-portal-impl}
% =====================================================================

Una volta autenticato, l'utente accede al portale (Fig.~\ref{fig:portal_menu}), che funge da punto di ingresso
unico per tutti gli agenti, in coerenza con il principio
di modularità verticale e accesso centralizzato
(Sez.~\ref{sec:agenti-principi}).

\begin{figure}[htbp]
  \centering
  \includegraphics[width=0.9\textwidth,height=0.5\textheight,keepaspectratio]{portale}
  \caption[Menu di scelta degli agenti nel portale]{Il portale utente dopo il login: ogni agente è raggiungibile con un click senza ulteriori autenticazioni.}
  \label{fig:portal_menu}
\end{figure}
\FloatBarrier

% =====================================================================
\section{Vera --- assistente generale con Web grounding}
\label{sec:agenti-vera-arch}
% =====================================================================
\label{sec:agenti-vera-impl}

Vera è l'agente generalista della suite e il più articolato dal punto di
vista realizzativo: questa sezione ne presenta il caso d'uso e le scelte
implementative significative, che ruotano attorno a tre problemi:
\emph{routing} del modello, \emph{Web grounding} con citazioni inline,
gestione dei documenti caricati dall'utente.

\subsection{Caso d'uso e requisiti}
\label{subsec:vera-caso-uso}

Vera deve coprire i casi d'uso tipici di un LLM generalista (analisi e riassunto di documenti,
formazione, correzione di codice e formule e così via): supporta il caricamento di documenti via
RAG~\cite{lewis2020rag} nel maggior numero di formati e, per le richieste che richiedono
informazioni attuali, accede a internet via Web grounding~\cite{gcp_grounding}, ancorando la
generazione ai risultati di Google Search. La fonte, da documento o dal web, deve essere sempre
presente e verificabile, e va garantito che la ricerca non esfiltri i dati --- potenzialmente
confidenziali --- presenti nei documenti o nei prompt dell'utente.

\subsection{Routing del modello basato sulla complessità della query}
\label{subsec:vera-routing}

La selezione del modello è governata dal modulo \texttt{model\_router.py},
che stima la complessità della richiesta invocando Gemini~2.5 Flash con un
prompt di classificazione dedicato (\texttt{thinking\_budget = 0} per
minimizzare la latenza). La risposta è vincolata a una singola parola,
\texttt{low} o \texttt{high}, e determina se il turno sarà servito da
Flash (bassa complessità) o da Gemini~2.5 Pro (alta). Lo stesso schema di
classificazione è adottato da tutti gli agenti della suite (SysAid,
AllyCare, CooPolicy), garantendo un trade-off costo/qualità uniforme come
descritto nella Sez.~\ref{sec:bg-routing}.

\subsection{Pattern di esecuzione: documenti vs Web grounding}
\label{subsec:vera-paths}

Il generatore distingue due percorsi mutuamente esclusivi sulla base della
presenza di documenti caricati nella conversazione. Quando ci sono \emph{upload}
in stato \texttt{ready}, viene costruita o riusata una \emph{Gemini Context
Cache} associata alla conversazione, che evita di ricaricare i medesimi
documenti a ogni turno; la generazione avviene quindi tramite l'agente ADK
\texttt{doc\_qa\_agent}, configurato in modo da utilizzare la cache
(\texttt{generate\_content\_config.cached\_content}).
In assenza di upload, il generatore percorre il \emph{path} di conversazione
generale, che monta un \texttt{LlmAgent} ADK con
l'\texttt{EnterpriseWebSearchTool} integrato: il modello decide
autonomamente, in base alla domanda, se invocare la \emph{search} o
rispondere con la sola conoscenza parametrica.

In entrambi i percorsi la frammentazione dell'output è gestita con la modalità
\texttt{StreamingMode.SSE} di ADK: il \emph{loop} di eventi del Runner emette
\texttt{Event} parziali, di cui vengono propagate solo le \texttt{parts} con
campo \texttt{partial = True}, escludendo i blocchi marcati come
\texttt{thought}. Nel percorso di conversazione generale, quando il modello
restituisce \texttt{finish\_reason} uguale a \texttt{MAX\_TOKENS}, il backend
invia automaticamente un turno addizionale ``Continua la risposta dal punto
esatto in cui si è interrotta'' fino a una soglia di sicurezza
(\texttt{\_MAX\_CONTINUATIONS = 4}), così da rendere la risposta completa
anche quando supera il limite di output del modello.

\subsection{Citazioni inline e \emph{grounding metadata}}
\label{subsec:vera-citations}

Quando il modello invoca la \emph{Web search}, l'evento ADK porta in un campo
dedicato (\texttt{event.grounding\_metadata}) la lista dei chunk consultati,
con URL, titolo e \emph{snippet}. La funzione \texttt{extract\_web\_citations}
trasforma questa struttura in un array di citazioni che il frontend rende
inline, cliccabili dall'utente. Per i documenti caricati la citazione viene
estratta da un \texttt{citation\_extractor} che attraversa il testo finale e
risolve i riferimenti contro la lista dei \texttt{cache\_uploads}. Il fatto
che la fonte sia sempre verificabile e cliccabile, e che la sua presenza sia
imposta dal sistema (non opzionale per il modello), realizza il vincolo di
verificabilità descritto nella Sez.~\ref{subsec:vera-caso-uso}.

Le Fig.~\ref{fig:vera_conv_web} e~\ref{fig:vera_conv_doc} mostrano due conversazioni reali con Vera, rispettivamente in modalità Web grounding (con fonti citate inline) e dopo il caricamento di un documento.

\begin{figure}[htbp]
  \centering
  \includegraphics[width=0.68\textwidth,height=0.85\textheight,keepaspectratio]{vera1}
  \caption[Conversazione con Vera: Web grounding]{Una conversazione con Vera in modalità Web grounding: le fonti recuperate da Google Search sono citate inline e cliccabili dall'utente.}
  \label{fig:vera_conv_web}
\end{figure}

\begin{figure}[htbp]
  \centering
  \includegraphics[width=\textwidth,keepaspectratio]{vera2}
  \caption[Conversazione con Vera: documento caricato]{Una conversazione con Vera dopo il caricamento di un documento: la risposta è ancorata al file tramite Gemini Context Cache e la fonte è indicata.}
  \label{fig:vera_conv_doc}
\end{figure}

\subsection{Catena di sicurezza pre e post generazione}
\label{subsec:agenti-security}

Coerentemente con il principio di governance facilitata, ogni invocazione di
modello in Vera attraversa una catena di controlli predisposta nel servizio
\texttt{model\_armor\_service.py}. La catena ha tre stadi:

\begin{enumerate}
  \item \textbf{Sanitizzazione del prompt.} Prima dell'invio all'LLM, il testo passa per il template
        di Google Cloud Model Armor~\cite{gcp_modelarmor} (Sez.~\ref{sec:bg-security}), che combina un classificatore
        \emph{Responsible AI} (\emph{hate}, \emph{harassment}, \emph{sexually explicit},
        \emph{dangerous}), il rilevatore di \emph{prompt injection}~\cite{perez2022promptinjection}
        e la \emph{Sensitive Data Protection}~\cite{gcp_dlp} per le PII, in linea con l'European AI
        Act~\cite{eu2024aiact}. Le PII del contesto italiano (Codice Fiscale, Partita IVA, IBAN) sono
        \emph{custom info types}; quando rilevate, il backend non maschera ciecamente il valore ma
        rifiuta la richiesta invitando a riformulare --- il \emph{tokenizing} cieco produrrebbe
        risposte inutilizzabili, poiché l'autenticità del dato è spesso necessaria.
  \item \textbf{Esenzioni per gruppo.} Alcuni gruppi (es. l'ufficio legale) hanno legittimo bisogno
        di trattare PII specifiche. Le tabelle \texttt{model\_armor\_group\_exemptions} e
        \texttt{pii\_group\_exemptions} consentono di escludere per gruppo singoli filtri o tipi di
        PII: alla verifica, il backend interseca le esenzioni attive dell'utente con i \emph{reasons}
        di Model Armor e, se ogni motivo è coperto, la richiesta prosegue.
  \item \textbf{Sanitizzazione della risposta e audit.} La stessa sanitizzazione si applica
        all'output prima del salvataggio, e ogni interazione è registrata in \texttt{chat\_logs}
        (action, modello, token in/out, flag citazioni). Gli eventi di blocco finiscono inoltre in un
        \emph{security log} dedicato che alimenta la dashboard di amministrazione
        (Sez.~\ref{sec:agenti-admin-impl}).
\end{enumerate}

\begin{figure}[htbp]
  \centering
  \includegraphics[width=0.85\textwidth,height=0.75\textheight,keepaspectratio]{flusso-prompt-full}
  \caption[Catena di sicurezza pre e post generazione]{Il flusso completo di una richiesta in Vera: verifica della quota, sanitizzazione del prompt tramite Model Armor
  (RAI + prompt injection + DLP) con gestione delle esenzioni per gruppo, chiamata al modello, sanitizzazione della risposta e
  audit su \texttt{chat\_logs}. Gli agenti verticali esclusivamente grounded non implementano questa catena (cfr.\ paragrafo successivo).}
  \label{fig:security_chain}
\end{figure}
\FloatBarrier

\paragraph{Ambito di applicazione.}
La catena è specifica di Vera; Orion (Cap.~\ref{chap:orion}) implementa la stessa, sanitizzando il
prompt prima di inoltrarlo ai sub-agenti. Gli altri agenti (SysAid, AllyCare, CooPolicy) non hanno
una propria catena pre/post generazione: in modalità \emph{standalone} operano solo sulla rete
interna con accesso autenticato (Sez.~\ref{subsec:agenti-auth}) e si affidano, come protezione a
livello di prompt, ai guardrail del proprio prompt di sistema (Sez.~\ref{sec:agenti-principi}).

% =====================================================================
\section{CooPolicy --- RAG su corpus di policy condiviso}
\label{sec:agenti-coopolicy-arch}
% =====================================================================
\label{sec:agenti-coopolicy-impl}

CooPolicy è la realizzazione ``RAG-native'' della suite: tutto il backend
ruota attorno alla pipeline di Retrieval-Augmented
Generation~\cite{lewis2020rag}, con scelte mirate a comprimere il rumore di
recupero e ad ancorare la generazione esclusivamente al materiale recuperato.

\subsection{Caso d'uso e requisiti}
\label{subsec:coopolicy-caso-uso}

L'agente è specializzato nella ricerca e interpretazione dei documenti normativi interni: gestisce
la documentazione della Intranet aziendale --- manuali, procedure, contratti di lavoro, organigrammi,
comunicazioni --- espansasi negli anni fino a diventare frastagliata e di non facile consultazione.

A differenza di Vera, CooPolicy non accede a fonti esterne: le risposte devono essere precise e
approfondite se il materiale è nel contesto RAG~\cite{lewis2020rag,guu2020realm}, oppure non essere
fornite affatto, sempre con fonti citate e accessibili. Per instradare anche richieste stringate, il
sistema classifica ogni query verso il corpus tematico più pertinente e applica una soglia di
rilevanza permissiva, delegando al modello la decisione finale sull'utilità dei chunk recuperati.

\subsection{Multi-corpus e classificazione della query}
\label{subsec:coopolicy-multicorpus}

Una prima scelta architetturalmente significativa è la rinuncia al
\emph{single corpus}. La documentazione interna è ripartita in dieci corpus
Vertex AI~\cite{gcp_vertex_rag} tematici (\texttt{policy-hr},
\texttt{contratti-smartworking}, \texttt{amministrazione-personale},
\texttt{procedure-aziendali}, \texttt{processi-regolamenti},
\texttt{governance-identita}, \texttt{sicurezza},
\texttt{punti-vendita-operativita}, \texttt{materiale-pdv},
\texttt{comunicazioni-generali}), ciascuno con un proprio descrittore in
italiano del dominio coperto. La motivazione è duplice: da un lato la
qualità del retrieval cala rapidamente al crescere della cardinalità dei
chunk indicizzati nello stesso spazio vettoriale; dall'altro, alcune
classi di documenti hanno regole di interpretazione
specifiche (i contratti collettivi vincolano la risposta in modo diverso
dalle policy aziendali) e mantenerle separate facilita l'evoluzione
indipendente e rende meno probabili ambiguità e risposte contraddittorie
o illogiche.

A monte del retrieval, una funzione \texttt{classify\_query} chiede a
Gemini~2.5 Flash di scegliere il \emph{corpus primario} ed eventualmente
un \emph{corpus secondario} fra quelli configurati, restituendo
strettamente un oggetto JSON con due chiavi. La scelta di Gemini Flash è
motivata dalla coppia di richieste essenziale di questo passo: bassa
latenza (sub-secondo) e basso costo unitario, condizione necessaria
perché il classificatore venga invocato a ogni richiesta senza pesare
sull'economia del sistema. Il prompt include il descrittore di ciascun
corpus disponibile.

\subsection{Retrieval e soglia di rilevanza}
\label{subsec:coopolicy-retrieval}

Il retrieval è eseguito in parallelo sui corpus selezionati attraverso
\texttt{rag.retrieval\_query} di Vertex AI RAG Engine~\cite{gcp_vertex_rag},
con \texttt{top\_k = 60} sul corpus primario e \texttt{top\_k = 30} sul
secondario quando presente. La ricerca è vettoriale densa: ogni query
viene trasformata in un embedding e confrontata per similarità coseno
con i vettori dei chunk indicizzati.

I chunk restituiti sono filtrati con una soglia
\texttt{RAG\_SCORE\_THRESHOLD~=~0.15} sulla distanza coseno: il valore è
permissivo, in coerenza con la decisione di delegare al
modello generativo (vincolato dal \emph{system prompt}) la decisione su
quali chunk siano effettivamente utili. L'aspettativa è che un modello in
modalità a bassa temperatura (Sez.~\ref{subsec:coopolicy-generation})
sappia rifiutare contesto irrilevante, mentre una soglia troppo restrittiva
rischia di escludere materiale marginalmente rilevante che però contiene
la risposta. I risultati di primario e secondario vengono uniti,
deduplicati per testo identico, ordinati per distanza coseno crescente
e troncati a \texttt{RAG\_TOP\_K} chunk totali.

\subsection{Arricchimento dei chunk e gestione dei metadati}
\label{subsec:coopolicy-meta}

Ogni PDF nel bucket GCS è accompagnato da un sidecar \texttt{.meta.json}
generato in fase di ingestione, contenente titolo leggibile e data del
documento. La funzione \texttt{\_enrich\_chunks\_with\_meta} scarica in
parallelo i sidecar dei chunk recuperati, con cache in-memory a TTL
un'ora, e arricchisce ciascun chunk con il titolo (dato che molti documenti avevano
nativamente nomi in esadecimale senza senso compiuto) e con la
data; quest'ultima è poi utilizzata dal modello per privilegiare la
fonte più recente in caso di ridondanza informativa. I file di metadati
(\texttt{.json}) e i tabulari grezzi (\texttt{.csv}), che possono
finire indicizzati per altri motivi, sono filtrati esplicitamente prima
di costruire il prompt, evitando che il modello citi rumore.

\subsection{Generazione vincolata e citazioni}
\label{subsec:coopolicy-generation}

Il prompt di sistema impone, in italiano, otto regole non negoziabili:
analisi integrale dei documenti recuperati prima di rispondere; risposta
basata esclusivamente sul materiale fornito; obbligatoriet\`a delle
citazioni \texttt{[N]} inline; massima precisione su numeri, articoli e
condizioni; divieto esplicito di parafrasi libere; tono professionale
italiano; consapevolezza temporale (preferire la fonte più recente,
avvisare se il documento ha più di due anni). Il modello è selezionato con
lo stesso classificatore di complessità descritto per Vera
(Sez.~\ref{subsec:vera-routing}): Gemini~2.5 Flash per richieste semplici,
Gemini~2.5 Pro per quelle complesse,
con \texttt{temperature = 0.1} e \texttt{max\_output\_tokens = 8192}. La
temperatura quasi zero massimizza l'aderenza al prompt e riduce
ulteriormente il rischio di allucinazione~\cite{xu2024hallucination}, in
linea con la funzione di un agente di tipo ``ricerca verificabile''.

Il post-processing della risposta è il dispositivo che garantisce
l'integrità delle citazioni. Le citazioni multiple
(\texttt{[1, 2, 3]}) vengono espanse in citazioni singole
(\texttt{[1][2][3]}); le citazioni con indice fuori dal range dei
documenti realmente forniti, se presenti,
vengono rimosse dal testo finale; gli indici sono infine rinumerati
in ordine di prima comparsa, in modo che la lista delle fonti
restituita al frontend rispecchi l'ordine di lettura. Per ciascuna
citazione viene anche estratto uno \emph{snippet} di 300~caratteri
centrato sulle keyword della query, utile per la \emph{preview} nel
frontend.

\subsection{ACL e proxy autenticato per i documenti}
\label{subsec:coopolicy-acl}

Il controllo di accesso avviene a monte, nel frontend orchestrato da
Orion: l'agente \texttt{coopolicy} è incluso nella lista degli agenti
disponibili solo per gli utenti che appartengono ad almeno un gruppo
Entra~ID presente nella \emph{whitelist} \texttt{DOCS\_ALLOWED\_GROUPS}.
Gli utenti non autorizzati non vedono l'agente e non possono quindi
raggiungerlo, eliminando la necessità di verifiche ridondanti a livello
di singolo endpoint. L'accesso al PDF originale, infine, non
avviene tramite URL pubblico: l'endpoint \texttt{/api/docs/file} agisce
come proxy autenticato che, verificato il token, ottiene un \emph{access
token} cloud-platform tramite le credenziali del service account del
servizio e scarica il binario direttamente dall'API di GCS, restituendolo
con \texttt{Content-Disposition: inline}. Il modello non ha mai
visibilit\`a degli URL firmati e l'utente non pu\`o aggirare la
verifica del token visitando un link.

Le Fig.~\ref{fig:coopolicy_conv1} e~\ref{fig:coopolicy_conv2} mostrano una conversazione con CooPolicy e il relativo pannello delle fonti, da cui il PDF originale è raggiungibile attraverso il proxy autenticato appena descritto.

\begin{figure}[htbp]
  \centering
  \includegraphics[width=0.68\textwidth,height=0.85\textheight,keepaspectratio]{coopolicy2}
  \caption[Conversazione con CooPolicy: risposta da corpus policy]{Una conversazione con CooPolicy: la risposta è basata esclusivamente sui documenti del corpus RAG,
  con citazioni numerate \texttt{[N]} che rimandano alla fonte originale.}
  \label{fig:coopolicy_conv1}
\end{figure}

\begin{figure}[htbp]
  \centering
  \includegraphics[width=0.95\textwidth,keepaspectratio]{coopolicy3}
  \caption[Pannello delle fonti di CooPolicy]{Il pannello delle fonti restituito da CooPolicy: ogni citazione numerata rimanda al documento originale,
  accessibile tramite il pulsante ``Apri PDF'' servito dal proxy autenticato verso GCS (Sez.~\ref{subsec:coopolicy-acl}).}
  \label{fig:coopolicy_conv2}
\end{figure}

% =====================================================================
\section{AllyCare / NPS --- Text-to-SQL sui dati dei sondaggi}
\label{sec:agenti-allycare-arch}
% =====================================================================
\label{sec:agenti-allycare-impl}

AllyCare è la realizzazione del paradigma Text-to-SQL
(Sez.~\ref{subsec:bg-text2sql}) su BigQuery per l'analisi dei sondaggi NPS.
L'implementazione espone al modello tre \emph{tool} di interazione con i
dati e affida all'\emph{agentic loop} di ADK la gestione del ciclo ReAct
(Sez.~\ref{sec:bg-multiagent}).

\subsection{Caso d'uso e requisiti}
\label{subsec:allycare-caso-uso}

L'agente è specializzato nell'analisi dei sondaggi NPS (Net Promoter Score) dell'organizzazione. La
base dati tabellare, in aggiornamento costante, raccoglie per ogni sondaggio i voti sui vari aspetti
dell'esperienza d'acquisto più un campo di testo libero; un flusso LLM dedicato analizza i commenti
e li cataloga (reparto, argomento, sentiment), producendo una tabella ``ai enriched'' che affianca ai
dati di base i tag generati.

AllyCare facilita le analisi su questa tabella, coprendo sia compiti da reportistica ad hoc (domande
quantitative come ``voti positivi nel periodo X per il negozio Y'', grafici di andamento) sia analisi
qualitative onerose (``elenca le maggiori criticità in una regione/negozio''). La richiesta è tradotta
dal modello in una o più query SQL, eseguite e usate per la risposta e le analisi successive.

\subsection{Iniezione selettiva dello schema}
\label{subsec:allycare-schema}

Il prompt di sistema include lo schema BigQuery delle due tabelle
rilevanti (\texttt{mapp\_nps\_survey\_internal\_prod} con i dati di
sondaggio grezzi e \texttt{mapp\_nps\_survey\_ai\_enriched\_prod} con i
campi prodotti da una pipeline di arricchimento AI a monte) ottenuto via
\texttt{bq\_client.get\_table} e mantenuto in cache per un'ora. Allo
schema tecnico viene affiancato un \emph{glossario semantico} in
italiano: per ciascuna colonna è descritto il significato funzionale
(es. \texttt{voto\_item01: ``Il personale del negozio è gentile e
competente''; valori -2 = mancata risposta, -1 = non lo so, 0--9 = voto}).
Insieme allo schema viene iniettata anche la lista esatta delle
descrizioni di negozio (\texttt{store\_desc}, in cache 24~ore) tratta
direttamente da BigQuery, in modo che il modello possa risolvere
varianti colloquiali (``Sasso'', ``Iper Bologna'') in match esatti senza
allucinare codici. Questa strategia di \emph{schema injection
selettiva}, in luogo dell'esposizione del catalogo completo del data
warehouse, è coerente con la letteratura
Text-to-SQL~\cite{pourreza2023dinsql}: meno colonne il modello vede,
meno opportunità ha di confondersi.

Il prompt impone inoltre regole di consistenza SQL specifiche del
dominio: i filtri sui voto-item devono includere \texttt{voto\_itemXX
>= 0} nella \texttt{WHERE} (escludendo i marcatori \texttt{-1} e
\texttt{-2}); tutti i confronti su campi categorici testuali devono
usare \texttt{UPPER(TRIM(...))} per neutralizzare incoerenze di
\emph{case}; la distinzione fra \texttt{LIKE} e uguaglianza esatta
sui nomi di negozio segue regole esplicite. L'effetto è che AllyCare
tipicamente non sbaglia i filtri ``stupidi'' che dominano gli errori in
Text-to-SQL applicato in dominio reale.

\subsection{I \emph{tool} disponibili l'\emph{agentic loop} ADK}
\label{subsec:allycare-tools}

Lo strato di tool è dichiarato in \texttt{tools.py} ed esporta tre
funzioni Python registrate come \emph{ADK tools}:

\begin{description}
  \item[\texttt{run\_sql(query)}]
    Esegue una query SQL su BigQuery con un \emph{guard} sintattico: la
    query è parsificata con \texttt{sqlparse} e respinta se contiene
    \texttt{DDL} o \texttt{DML} di scrittura
    (\texttt{INSERT}/\texttt{UPDATE}/\texttt{DELETE}); se non già
    presente, viene aggiunto un \texttt{LIMIT~5000}; il job è eseguito
    con \texttt{maximum\_bytes\_billed = 10\textsuperscript{9}} per
    bloccare query accidentalmente costose. Il risultato è restituito
    al modello come tabella Markdown delle prime 750~righe, con
    indicazione del totale; la copia completa del DataFrame è
    memorizzata nello \texttt{tool\_context.state} per essere
    riutilizzata da \texttt{generate\_chart} senza una seconda
    interrogazione. Se la query produce un errore SQL, il messaggio
    viene impacchettato come \texttt{function\_response} e il modello
    può auto-correggersi al ciclo successivo: questo è il meccanismo
    concreto di \emph{self-correction}~\cite{pourreza2023dinsql}.
  \item[\texttt{analyze\_comments(query, question)}]
    Estrae un campione di commenti testuali grezzi mediante una query
    dedicata, ne preleva un sotto-campione casuale di al più 200 record
    e chiede a Gemini~2.5 Pro una sintesi qualitativa con citazioni
    \emph{verbatim}. Il prompt di analisi è scritto in modo da essere
    integrato nella risposta più ampia e da non duplicare l'intestazione
    di sezione. È usato esclusivamente per analisi qualitative: per
    qualunque dato quantitativo (sentiment, conteggi di temi, percentuali)
    il prompt di sistema impone \texttt{run\_sql} sulla tabella
    \texttt{ai\_enriched\_prod}.
  \item[\texttt{generate\_chart(visualization\_goal)}]
    Costruisce un grafico Matplotlib/Seaborn a partire dal DataFrame
    cachato dall'ultimo \texttt{run\_sql}, scegliendo la rappresentazione
    in funzione di parole-chiave nel \emph{goal}
    (\texttt{torta}/\texttt{pie} $\rightarrow$ pie chart;
    \texttt{linee}/\texttt{trend} $\rightarrow$ line chart; default
    $\rightarrow$ bar chart). Il grafico è codificato in PNG, convertito
    in base64 con prefisso \emph{data URI} e iniettato nel session state.
    Il backend lo restituisce poi nella risposta come campo
    \texttt{chart\_and\_data}, lasciando il rendering alla logica del
    frontend. Il chart viene incluso nel payload finale solo se
    \texttt{generate\_chart} è stato effettivamente invocato nel turno,
    evitando di restituire il grafico del turno precedente.
\end{description}

L'\emph{agentic loop} è governato dal \texttt{Runner} ADK che, a partire
dall'\texttt{LlmAgent} configurato con i tre tool e con il prompt di
sistema, alterna in modo trasparente fasi di ragionamento e chiamate a
tool fino a quando il modello restituisce la risposta finale. Il
modello scelto per il turno (Pro o Flash) è iniettato nell'agente al
momento della costruzione, e il prompt di sistema include il nome
proprio dell'utente e la sua direzione (estratti dal token Entra ID),
così che il tono della risposta sia personalizzato.

La selezione del modello segue lo stesso schema adottato per SysAid
(Sez.~\ref{subsec:sysaid-arch-unificata}): il campo \texttt{complexity}
propagato da Orion determina se il turno è servito da Gemini~2.5 Flash
(\texttt{low}) o Gemini~2.5 Pro (\texttt{high}). Quando AllyCare è
usato in modalità standalone, la complessità è stimata localmente dal
modello Flash con una chiamata di classificazione dedicata,
coerentemente con il pattern già descritto per SysAid.

\subsection{Disciplina del ragionamento, anti-allucinazione}
\label{subsec:allycare-discipline}

La ``gerarchia della verità'' codificata nel prompt impone che i dati
del database siano la sola fonte autoritativa: il modello non può
classificare commenti per produrre numeri (la categorizzazione
quantitativa va sempre dalla tabella arricchita). Una regola di
``fiducia nel tool'' impone inoltre che, se \texttt{run\_sql} ha
restituito un dato, sia quel dato a essere riportato, anche se il
modello ``ricorda'' numeri diversi da turni precedenti.

Le Fig.~\ref{fig:allycare_conv1}, \ref{fig:allycare_chart} e~\ref{fig:allycare_conv2} mostrano tre interazioni con AllyCare: un'analisi quantitativa dell'NPS, un grafico generato dinamicamente dai dati estratti via SQL e un'analisi qualitativa dei commenti liberi.

\begin{figure}[htbp]
  \centering
  \includegraphics[width=0.68\textwidth,height=0.85\textheight,keepaspectratio]{nps2}
  \caption[Conversazione con AllyCare: analisi quantitativa]{Una conversazione con AllyCare: la richiesta di analisi NPS produce una sezione di \emph{ragionamento}
  esplicito (con la strategia SQL adottata) seguita da una risposta numerica strutturata (punteggio NPS, suddivisione promotori/passivi/detrattori).}
  \label{fig:allycare_conv1}
\end{figure}

\begin{figure}[htbp]
  \centering
  \includegraphics[width=0.6\textwidth,height=0.8\textheight,keepaspectratio]{nps5}
  \caption[AllyCare: grafico generato dinamicamente]{AllyCare genera dinamicamente un grafico (qui un grafico a torta della distribuzione del \emph{sentiment} dei commenti) a
  partire dai dati estratti via SQL, codificandolo come immagine restituita al frontend.}
  \label{fig:allycare_chart}
\end{figure}

\begin{figure}[htbp]
  \centering
  \includegraphics[width=0.68\textwidth,height=0.85\textheight,keepaspectratio]{nps4}
  \caption[Conversazione con AllyCare: analisi qualitativa dei commenti]{AllyCare in modalità analisi qualitativa: i commenti testuali liberi del sondaggio sono sintetizzati per tema,
  con la relativa distribuzione di \emph{sentiment}, senza produrre numeri non supportati dalla tabella arricchita.}
  \label{fig:allycare_conv2}
\end{figure}

% =====================================================================
\section{SysAid --- Text-to-SQL sul ticketing con modelli LLM distinti}
\label{sec:agenti-sysaid-arch}
% =====================================================================
\label{sec:agenti-sysaid-impl}

SysAid è l'agente in cui il principio di indipendenza dal fornitore di LLM
trova la sua realizzazione più esplicita, con un backend unificato che
ospita due code path alimentate da fornitori concorrenti.

\subsection{Caso d'uso e requisiti}
\label{subsec:sysaid-caso-uso}

SysAid è specializzato nell'analisi della base dati di ticketing dell'organizzazione, una copia
tabellare in streaming in tempo reale del database del servizio. Come AllyCare, copre sia compiti da
reportistica ad hoc (numeriche, andamenti, tempi di risoluzione) sia analisi qualitative onerose
(``maggiori criticità in una regione/negozio''); la disponibilità di dati realtime gli consente
inoltre di riportare lo stato di salute attuale dei sistemi a supporto del troubleshooting --- utile
al management e al personale del Service Desk. Anche qui la richiesta è tradotta in query SQL,
eseguite e usate per la risposta.

\subsection{Pluralità del modello nel routing}
\label{subsec:agenti-sysaid-confronto}

Il backend deve essere realizzato come servizio unificato che ospita due code parallele, alimentate da due fornitori concorrenti:
Gemini 2.5 Flash su Vertex AI~\cite{gcp_vertexai} per le richieste a bassa complessità e
Claude Sonnet su AnthropicVertex~\cite{anthropic_docs} per le richieste ad alta complessità;
Gemini Flash si deve occupare anche della valutazione della complessità.

Entrambe le code path condividono la medesima base di conoscenza, la stessa struttura di database,
i medesimi strumenti per l'analisi dei dati e l'interfaccia utente.

\subsection{Architettura del backend unificato}
\label{subsec:sysaid-arch-unificata}

Il backend unificato condivide fra le due code path tutti gli elementi non
legati al fornitore del modello: schema BigQuery (vista \texttt{service\_req\_view}
sulla base ticket SysAid), prompt di sistema, set di tool
(\texttt{run\_sql}, \texttt{analyze\_text}, \texttt{generate\_chart}),
conversazioni e interfaccia utente React. Il routing interno avviene
all'ingresso dell'endpoint \texttt{/api/sysaid/prompt}: il campo
\texttt{complexity} determina quale code path viene percorso. La
Tab.~\ref{tab:sysaid-confronto} confronta le due code path, evidenziando come la logica di business resti condivisa e le differenze si concentrino sul solo fornitore di LLM.

\begin{table}[htbp]
  \centering
  \caption[Code path nel backend SysAid unificato]{Confronto fra le due
    code path ospitate nel backend SysAid unificato. La logica di business
    è condivisa; le differenze riguardano esclusivamente il fornitore di LLM.}
  \label{tab:sysaid-confronto}
  \begin{tabular}{lll}
    \toprule
    \textbf{Dimensione}      & \textbf{Code path Gemini (low)} & \textbf{Code path Sonnet (high)}  \\
    \midrule
    Modello LLM              & Gemini 2.5 Flash                & Claude Sonnet 4.6                 \\
    Piattaforma              & Vertex AI (\texttt{europe-west8}) & AnthropicVertex (\texttt{europe-west1}) \\
    SDK                      & Google ADK nativo               & Anthropic SDK + thin adapter      \\
    Schema dei tool          & \texttt{FunctionDeclaration}    & JSON Schema                       \\
    Stop reason tool         & \texttt{function\_call} part    & \texttt{stop\_reason = tool\_use} \\
    Schema BigQuery          & \multicolumn{2}{c}{Condiviso}                                      \\
    Prompt di sistema        & \multicolumn{2}{c}{Identico}                                       \\
    Tool disponibili         & \multicolumn{2}{c}{\texttt{run\_sql}, \texttt{analyze\_text}\footnotemark, \texttt{generate\_chart}} \\
    Frontend React           & \multicolumn{2}{c}{Unico}                                          \\
    \bottomrule
  \end{tabular}
\end{table}
\footnotetext{L'interfaccia dei tool è condivisa, ma \texttt{analyze\_text}
  utilizza internamente il modello LLM della rispettiva code path: Gemini~2.5
  Flash nel percorso Gemini, Claude Sonnet~4.6 nel percorso Sonnet.}

Una nota architetturale è la differenza di regione GCP fra le
due code path: Gemini è disponibile in \texttt{europe-west8}, mentre
Claude Sonnet su Vertex è offerto al momento solo in
\texttt{europe-west1}. Entrambe le regioni si trovano nel perimetro
europeo e soddisfano il requisito di sovranità del dato.

\subsection{Code path Gemini: integrazione ADK nativa}
\label{subsec:sysaid-gemini}

La code path Gemini utilizza appieno il \emph{Google Agent Development Kit} (ADK)
per l'astrazione del modello e l'orchestrazione del \emph{function calling}. I tool sono
registrati nativamente tramite i decoratori e le classi fornite dall'SDK,
delegando all'\emph{agentic loop} di ADK la gestione del ciclo Reason-Act-Observe.
Il modello è configurato con \texttt{temperature = 0},
che, in continuità con CooPolicy e con la letteratura sulla
calibrazione~\cite{kadavath2022knowsknows}, riduce la varianza
dell'output e rende riproducibile il comportamento su query
identiche.

\subsection{Code path Sonnet: astrazione tramite adapter}
\label{subsec:sysaid-sonnet}

La code path Sonnet accede a Claude Sonnet
instradando le chiamate verso \texttt{AnthropicVertex},
l'integrazione che permette di chiamare i modelli Anthropic mantenendo l'autenticazione, i dati
e la fatturazione all'interno del progetto GCP~\cite{anthropic_docs}. Poiché le API di Anthropic
e i formati di tool calling differiscono parzialmente dallo standard Gemini,
l'implementazione fa uso di un \emph{thin adapter} per uniformare la semantica.

\paragraph{Definizione dei tool.} I tool rimangono concettualmente identici e
vengono traslati nel formato JSON~Schema richiesto da Claude.

\paragraph{Ciclo e \texttt{stop\_reason}.} L'agentic loop si adatta per interpretare
lo \texttt{stop\_reason = "tool\_use"} caratteristico di Anthropic.
L'esecuzione del loop rimane gestita centralmente, garantendo lo stesso limite di iterazioni
e lo stesso approccio strutturato visto nella code path Gemini. Quando il tool loop termina,
un \emph{guardrail} specifico verifica la corretta conformità del Markdown di output, forzando
se necessario un retry automatico del modello in caso di fallimenti semantici o assenza di esecuzione SQL.

\paragraph{Serializzazione dei content block.} L'SDK Anthropic
restituisce content block che includono campi interni
(\texttt{caller}, \texttt{parsed\_output}, \ldots) non accettati dall'API
quando il blocco viene rispedito come parte dell'history. Per evitare
errori 400, un helper applica una
\emph{whitelist} per tipo di blocco
(\texttt{text}, \texttt{tool\_use}, \texttt{tool\_result},
\texttt{image}) che mantiene solo i campi documentati.

\paragraph{Gestione di output lunghi.} Per
sostenere \texttt{max\_tokens = 32000} senza incorrere in timeout
dell'SDK, il backend usa internamente l'API \emph{streaming} di
Anthropic (\texttt{messages.stream}), accumulando il messaggio finale
prima di processarlo nel loop. Quando \texttt{stop\_reason = "max\_tokens"},
il backend aggiunge un messaggio utente ``Continua\ldots'' e ripete il giro, in
analogia con il meccanismo di continuation di Vera (Sez.~\ref{subsec:vera-paths}).

L'effetto netto dell'architettura unificata è che la logica di business
(esecuzione SQL, analisi testuale, grafici, persistenza) è scritta una sola volta;
gli adattamenti rispetto al fornitore restano localizzati alla traduzione
di schema dei tool, all'ispezione dei \texttt{stop\_reason} e alla
serializzazione dei content block. Aggiungere un terzo fornitore richiederebbe
un nuovo adapter, non una riscrittura della logica.

Le Fig.~\ref{fig:sysaid_conv_gemini} e~\ref{fig:sysaid_conv_sonnet} mostrano due conversazioni reali con SysAid, servite rispettivamente dalla code path Gemini~2.5 Flash (bassa complessità) e dalla code path Claude Sonnet~4.6 (alta complessità).

\begin{figure}[htbp]
  \centering
  \includegraphics[width=0.68\textwidth,height=0.85\textheight,keepaspectratio]{sysaid3}
  \caption[Conversazione SysAid: code path Gemini Flash]{SysAid con code path Gemini~2.5 Flash per una richiesta a bassa complessità:
  la query SQL è eseguita con bassa latenza e il risultato è qui reso come grafico a torta della distribuzione degli stati dei ticket.}
  \label{fig:sysaid_conv_gemini}
\end{figure}

\begin{figure}[htbp]
  \centering
  \includegraphics[width=0.68\textwidth,height=0.85\textheight,keepaspectratio]{sysaid2}
  \caption[Conversazione SysAid: code path Claude Sonnet]{SysAid con code path Claude Sonnet~4.6 per una richiesta ad alta complessità:
  la logica di business è identica alla code path Gemini; cambia solo il modello che esegue il ragionamento e la generazione SQL.}
  \label{fig:sysaid_conv_sonnet}
\end{figure}

% =====================================================================
\section{Pannello di amministrazione}
\label{sec:agenti-admin-impl}
% =====================================================================

Il monitoraggio e la configurazione della suite sono affidati a un pannello
di amministrazione dedicato a Vera (\texttt{admin-backend} +
\texttt{admin-frontend}),  isolato dagli applicativi
rivolti all'utente finale. Il backend FastAPI espone i propri router sotto
il prefisso \texttt{/api/admin/} e offre le seguenti funzionalità:
la gestione dei budget token per utente (con la tabella condivisa
\texttt{token\_budget} descritta sopra), le statistiche di consumo aggregate,
i log audit delle conversazioni con il feedback strutturato degli utenti,
la configurazione dei gruppi di esenzione PII e la gestione dei template
Model Armor e DLP (Sensitive Data Protection).

Il pannello di Orion, descritto nel
Cap.~\ref{chap:orion}, replica le stesse funzionalità
e aggiunge le pagine specifiche per la valutazione del routing.

La qualità di questi agenti (accuratezza delle query Text-to-SQL di
AllyCare e SysAid e fedeltà delle risposte RAG di CooPolicy) è valutata
nel Capitolo~\ref{chap:orion-valutazione}
(Sez.~\ref{sec:eval-agenti}).

% =====================================================================
\section{Alternative scartate}
\label{sec:agenti-alternative}
% =====================================================================

Questa sezione documenta le principali alternative considerate e le
ragioni per cui non sono state adottate.

\paragraph{Orchestratore centrale fin dalla versione~1.}
La prima alternativa valutata era introdurre un orchestratore
intelligente già nella versione iniziale della suite, ispirato
ai pattern multi-agente proposti dalla letteratura recente~\cite{wu2023autogen,guo2024multiagent,wang2024moa}, eliminando la
necessità per l'utente di selezionare esplicitamente l'agente. Questa
opzione è stata inizialmente scartata per una ragione metodologica: costruire
contemporaneamente gli agenti specializzati e il loro orchestratore
avrebbe reso molto più difficile isolare e diagnosticare i problemi e probabilmente avrebbe
condotto a performance degli agenti ``verticali'' molto meno ottimizzate.
Validare ogni agente indipendentemente è molto più semplice in un contesto ``verticale'' isolato.
L'orchestratore è stato quindi introdotto solo a valle della validazione
dei singoli agenti, ed è oggetto del Capitolo~\ref{chap:orion}.

\paragraph{Fine-tuning o training di un modello proprietario o open source.}
Si è valutato di addestrare un modello sui dati interni per creare un LLM specializzato sul dominio.
L'opzione è stata scartata per due ragioni. Primo, il costo di addestramento e mantenimento su scala
enterprise è proibitivo per un'organizzazione il cui core business non è lo sviluppo di modelli AI
(manca l'hardware adeguato e acquisirlo non sarebbe giustificabile). Secondo, il fine-tuning non
risolve il \emph{knowledge cutoff}: un modello addestrato sui dati di sei mesi fa non conosce i ticket
di questa settimana né le policy aggiornate ieri --- per i dati che cambiano nel tempo la risposta
corretta è il grounding.
